\section{Package Architecture and Zero-Dependency Core Isolation}

To maintain high execution velocity and reproducibility on minimal environments, \texttt{NeuroARC.jl} enforces zero-dependency core isolation. The core library relies exclusively on the Julia standard library, implementing lightweight internal parsers and memory buffers without external binary dependencies.

\subsection{Standard Julia Package Skeleton}
The repository adheres to the standard Julia package layout:
\begin{itemize}
  \item \texttt{Project.toml}: Declares package UUID\\
  (\texttt{a3f5a15a-93e1-4bd8-99e6-0563b78be05c}), version bounds on Julia (\texttt{1.10} to \texttt{1.13}), and conditional extension dependencies.
  \item \texttt{src/NeuroARC.jl}: Exposes public APIs for grid representations, component analysis, criterial calculus, and hyphal grafting contracts.
  \item \texttt{ext/}: Houses conditional package extensions that load only when external packages are present in the environment:
  \begin{itemize}
    \item \texttt{NeuroARCMetalExt.jl}: GPU acceleration for Apple Silicon via \texttt{Metal.jl}.
    \item \texttt{NeuroARCJACCExt.jl}: Cross-hardware portability via \texttt{JACC.jl}.
    \item \texttt{NeuroARCTuringExt.jl}: Bayesian MCMC calibration via \texttt{Turing.jl}.
  \end{itemize}
  \item \texttt{test/runtests.jl}: Comprehensive unit test suite executing 652 assertions across core types, criterial gates, constructive certificates, grafting contracts, and epistemic text cleanliness.
\end{itemize}

\subsection{Continuous Integration and Clean Runners}
The repository includes GitHub Actions CI workflows configured for clean runners on Ubuntu and macOS. These workflows verify that the package compiles and passes all unit tests without pre-installing non-standard dependencies, ensuring that clean installations remain fully functional.
